\documentclass[10pt,a4paper]{amsart}
\usepackage[margin=19mm]{geometry}
\usepackage{amsmath,amssymb,mathtools}
\usepackage[scr=rsfs,cal=euler]{mathalfa}
\usepackage{xfrac}
\usepackage[unicode,hidelinks,bookmarks=false]{hyperref}
\newcommand{\ie}{i.e.\ }
\newcommand{\cf}{cf.\ }
\newcommand{\resp}{respectively\ }
\newcommand{\Subsection}{\subsection}
\providecommand{\nobreakdash}{\nobreak\hbox{-}}
\setlength{\emergencystretch}{2em}
\multlinegap0pt
\theoremstyle{plain}
\newtheorem{theorem}{Theorem}
\newtheorem{lemma}[theorem]{Lemma}
\newtheorem{proposition}[theorem]{Proposition}
\newtheorem{corollary}[theorem]{Corollary}
\theoremstyle{definition}
\newtheorem{definition}[theorem]{Definition}
\theoremstyle{remark}
\newtheorem{remark}[theorem]{Remark}
\usepackage{bbm}
\newcommand{\noi}{\noindent}
\newcommand{\Z}{\mathbb{Z}}
\newcommand{\R}{\mathbb{R}}
\newcommand{\T}{\mathbb{T}}
\renewcommand{\H}{\mathcal H}
\newcommand{\E}{\mathbb{E}}
\newcommand{\dl}{\delta}
\newcommand{\nb}{\nabla}
\newcommand{\ep}{\varepsilon}
\newcommand{\g}{\gamma}
\newcommand{\cj}{\overline}
\newcommand{\les}{\lesssim}
\newcommand{\jb}[1]
{\langle #1 \rangle}
\newcommand{\ind}{\mathbbm 1}
\newcommand{\ph}{\varphi}
\newcommand{\eps}{\ep}
\let\Re=\undefined\DeclareMathOperator*{\Re}{Re}
\newcommand{\HF}[2]{\mathcal H\big(\tfrac1p{\textstyle\int_\T}|\cdot|^p\big)(#1)[#2]}
\newcommand{\WV}[1]{e^{\frac1p\int_\T|#1+D_t^{-1}\jb\nb u|^p-\frac1p\int_\T|#1|^p-\langle D_t^{-1}\jb\nb u,\,|#1|^{p-2}#1\rangle}\,\ind_{\{\|#1+D_t^{-1}\jb\nb u\|_{L^2}^2\le K\}}}
\newcommand{\RT}{R_v(D_t^{-1}\jb\nb T_{t,v}u)}
\newcommand{\IT}{\ind_{\{\|v+D_t^{-1}\jb\nb T_{t,v}u\|_{L^2}^2\le K\}}}
\title[A remark on the log-Sobolev inequality for the Gibbs measure of the focusing Schrodinger equation]{A remark on the log-Sobolev inequality for the Gibbs measure of the focusing Schrodinger equation: the spectral-gap criterion for the truncated Gibbs measure (printed as Proposition 4.6)}
\author{Guopeng Li, Jiawei Li, Leonardo Tolomeo}
\date{Source: arXiv:2512.03897v2 (CC BY 4.0) (CC BY 4.0)}
\begin{document}
\maketitle

\noindent\emph{Source and licence.} A remark on the log-Sobolev inequality for the Gibbs measure of the focusing Schrodinger equation. Version arXiv:2512.03897v2 (CC BY 4.0), distributed under the Creative Commons Attribution 4.0 International licence, http://creativecommons.org/licenses/by/4.0/, as recorded in the arXiv licence metadata for this exact version and shown on the saved abstract page https://arxiv.org/abs/2512.03897v2. Full licence notice: licenses/arxiv-2512.03897-CC-BY-4.0.txt.\par\smallskip

\noindent\textit{Editorial scope.} Counted unit: the quantitative spectral-gap criterion for the truncated Gibbs measure of the focusing nonlinear Schr\"odinger equation on the torus, printed in the article as Proposition 4.6 (source0.tex lines 1589--1610, one \texttt{proposition} environment with source label \texttt{s4:PROP:crit}), delivered together with the article's complete original proof of it (source0.tex lines 1612--1690, one \texttt{proof} environment of 79 lines immediately adjacent to the statement, with no gap and no deferral). The proof is the article's own quantitative argument: it bounds the auxiliary parameter by the exponential-moment estimate, combines the recentring lemma, Taylor's formula and the two-sided bound for the weighted Hessian to compare the Hessian of the renormalised potential with the Gaussian covariance term up to a factor controlled by the small parameter, evaluates the resulting ratio along the recentred field, and closes the estimate with the Gaussian lemma. No mathematics is written, repaired or re-derived: statement and proof are the article's own text line for line, in the article's own notation, and no line of either is omitted, substituted or re-flowed.

Required context retained uncounted: the display defining the second variation of the nonlinearity (lines 1344--1347); the two-sided bound for the weighted Hessian (lines 1351--1353); Taylor's formula with the remainder bound (lines 1356--1362); the recentring lemma with the display that defines the recentred field (lines 1366--1384); the Gaussian lemma with its two estimates (lines 1426--1458); and the exponential-moment lemma (lines 1506--1516). Those are exactly the labelled statements the proposition invokes, and nothing else from the article's development is used.

References. No citation command occurs in any retained line, so this package carries no bibliography block and no bibliography entry.

Editorial formatting only, in addition: an AMS \texttt{amsart} document, the same class the article itself uses, that restates verbatim the article's own macro definitions needed by the retained lines, loads the standard bbm package for the article's own indicator symbol, and restates the article's own theorem-environment declarations with the same shared counter the article declares. Consequently the retained environments print here in the order in which they occur, so the counted proposition prints as Proposition 4, whereas the article prints it as Proposition 4.6 under its per-section numbering; the retained lemmas print as Lemmas 1, 2 and 3 under the same shared counter, the equations of the retained spans are numbered anew in order of appearance, and every cross-reference inside the retained text resolves to this wrapper's numbering. That is the only change to the numbering. No figure, no image-inclusion line and no drawing code occur in any retained line, so no artwork is delivered or needed.

Not retained: the article's introduction, its two main theorems, its other lemmas and propositions, its appendix and its bibliography; none of that material is used as a step of the delivered proof, and the delivered proof is complete as the author wrote it. The package ships the article's concatenated source verbatim as \texttt{sources/190291/source0.tex}, and every delivered excerpt is an exact line range of that file. No mathematical statement, hypothesis, estimate or conclusion has been rewritten, repaired or added, and printed slips, if any, are preserved literally.
\begin{equation}\label{s4:secvar}
\HF{v}{g,g}:=\frac{d^2}{d\eps^2}\Big|_{\eps=0}\frac1p\int_\T|v+\eps g|^pdx
=\int_\T\Big(|v|^{p-2}|g|^2+(p-2)|v|^{p-4}\big(\Re\,\cj vg\big)^2\Big)dx ,
\end{equation}

\begin{equation} \label{HintUB}
    \int_\T|v|^{p-2}|g|^2dx\le\HF{v}{g,g}\le(p-1)\int_\T|v|^{p-2}|g|^2dx .
\end{equation}

\begin{equation}\label{s4:Taylor}
\begin{aligned}
&\frac1p\int_\T|v+\dl|^p-\frac1p\int_\T|v|^p-\big\langle\dl,|v|^{p-2}v\big\rangle
=\frac12\HF{v}{\dl,\dl}+R_v(\dl),\\
&|R_v(\dl)|\le C_p\int_\T\big(|v|^{p-3}|\dl|^3+|\dl|^p\big)dx .
\end{aligned}
\end{equation}

\begin{lemma}\label{s4:LEM:recentre}
Let $t\in(0,\infty)$ and $v\in H^1(\T)\cap L^\infty(\T)$, and set
\begin{equation}\label{s4:phiv}
\ph:=v-D_t^{-2}\big(|v|^{p-2}v\big)\in H^1(\T).
\end{equation}
Then the image of $\rho^t_\ph$ under the translation
$u\mapsto u-D_t^{-1}\jb\nb^{-1}\big(|v|^{p-2}v\big)$ is the probability measure
\begin{equation}\label{s4:rhotilde}
d\tilde\rho^{\,t}_v(u):=\frac1{\tilde Z^t_v}\,\WV{v}\,d\mu(u),
\end{equation}
where $\tilde Z^t_v$ is the normalisation constant. Consequently, for every
trigonometric polynomial $h$,
\begin{equation}\label{s4:Hessv}
\H V_t(\ph)[h,h]
=\E_{\tilde\rho^{\,t}_v}\big[\jb{u,D_t\jb\nb h}^2\big]-\Big(\E_{\tilde\rho^{\,t}_v}\big[\jb{u,D_t\jb\nb h}\big]\Big)^2-\|D_th\|_{L^2}^2 .
\end{equation}
By \eqref{s4:Taylor}, the exponent in \eqref{s4:rhotilde} equals
$\frac12\HF{v}{D_t^{-1}\jb\nb u,D_t^{-1}\jb\nb u}+R_v(D_t^{-1}\jb\nb u)$.
\end{lemma}

\begin{lemma}
\label{s4:LEM:gauss}
Let $t>0$ and $v\in L^\infty(\T)$, and let
\begin{equation}\label{s4:theta}
\theta_t(v):=\sup_{f\neq0}\frac{\HF{v}{D_t^{-1}f,D_t^{-1}f}}{\|f\|_{L^2}^2}
\end{equation}
be the norm of the non-negative self-adjoint operator $D_t^{-1}\mathfrak h_vD_t^{-1}$ on $L^2$. Then
\begin{equation}\label{s4:thetabound}
\theta_t(v)\le(p-1)\sum_{n\in\Z}\frac{1-e^{-t\jb n^2}}{\jb n^2}\int_\T|v|^{p-2}dx .
\end{equation}
Assume moreover that $\theta:=\theta_t(v)<1$, and let
\[
T_{t,v}:=\jb\nb^{-1}\big(1-D_t^{-1}\mathfrak h_vD_t^{-1}\big)^{-\frac12}\jb\nb .
\]
Then the following hold.
\begin{enumerate}
\item[(i)] For every measurable
$F\ge0$,
\[
\frac{\E_\mu\big[F(u)\,e^{\frac12\HF{v}{D_t^{-1}\jb\nb u,D_t^{-1}\jb\nb u}}\big]}{\E_\mu\big[e^{\frac12\HF{v}{D_t^{-1}\jb\nb u,D_t^{-1}\jb\nb u}}\big]}
=\E_\mu\big[F(T_{t,v}u)\big].
\]
\item[(ii)] $\|T_{t,v}g\|_{H^1}^2\le(1-\theta)^{-1}\|g\|_{H^1}^2$ for every $g$.
\item[(iii)] $\|D_t^{-1}\jb\nb T_{t,v}g\|_{H^1}^2\le(1-\theta)^{-1}\|g\|_{H^1}^2$ for every $g$.
\item[(iv)] For every $x\in\T$, $\displaystyle\E_\mu\big|D_t^{-1}\jb\nb T_{t,v}u(x)\big|^2\le\frac2{1-\theta}\sum_{n\in\Z}\frac{1-e^{-t\jb n^2}}{\jb n^2}$.
\item[(v)] For every trigonometric polynomial $h$,
\[
\E_\mu\big[\jb{T_{t,v}u,D_t\jb\nb h}^2\big]-\|D_th\|_{L^2}^2\ \ge\ \HF{v}{h,h},
\qquad
\E_\mu\big[\jb{T_{t,v}u,D_t\jb\nb h}^2\big]\le\frac{\|D_th\|_{L^2}^2}{1-\theta}.
\]
\end{enumerate}
\end{lemma}

\begin{lemma}[Exponential moments of the remainder]\label{s4:LEM:BD}
Let $t>0$, $q\ge1$ and $\theta_0\in(0,1)$, and let $v\in L^\infty(\T)$ with
\[
\|v\|_{L^2}^2\le K,\qquad (p-1)\sum_{n\in\Z}\frac{1-e^{-t\jb n^2}}{\jb n^2}\int_\T|v|^{p-2}dx\le\theta_0 .
\]
Then
\[
\log\E_\mu\Big[e^{q|R_v(D_t^{-1}\jb\nb T_{t,v}u)|}\ind_{\{\|v+D_t^{-1}\jb\nb T_{t,v}u\|_{L^2}^2\le K\}}\Big]\le C,
\]
where $C<\infty$ depends only on $p$, $q$, $\theta_0$ and $K$.
\end{lemma}

\begin{proposition}
\label{s4:PROP:crit}
Let $\theta_0\in(0,1)$. There exists $\gamma_0\ll1$,
depending only on $p$, $K$ and $\theta_0$, such that the following holds; the implicit constant in \eqref{s4:crit} below also depends only on $p$, $K$ and $\theta_0$. Let $t\in(0,1]$ and
$v\in H^1\cap L^\infty(\T)$ satisfy
\[
\|v\|_{L^2}^2\le\frac K2,\qquad (p-1)\sum_{n\in\Z}\frac{1-e^{-t\jb n^2}}{\jb n^2}\int_\T|v|^{p-2}\le\theta_0,
\qquad
\gamma:=t^{\frac34}\int_\T|v|^{p-3}+t^{\frac14}\le\gamma_0 .
\]
Let $\ph$ be given by \eqref{s4:phiv}. Then, for every trigonometric polynomial $h$,
\begin{equation}\label{s4:crit}
\Big|\H V_t(\ph)[h,h]-\Big(\E_\mu\big[\jb{T_{t,v}u,D_t\jb\nb h}^2\big]-\|D_th\|_{L^2}^2\Big)\Big|
\les \gamma\,\|D_th\|_{L^2}^2 .
\end{equation}
% In particular, for every trigonometric polynomial $w$,
% \begin{equation}\label{s4:critw}
% \begin{aligned}
% &\Big|\H V_t(\ph)\big[e^{-\frac t2\jb\nb^2}w,e^{-\frac t2\jb\nb^2}w\big]-\Big(\E_\mu\big[\jb{T_{t,v}u,D_t\jb\nb e^{-\frac t2\jb\nb^2}w}^2\big]-\|D_te^{-\frac t2\jb\nb^2}w\|_{L^2}^2\Big)\Big|\les \gamma\,t^{-1}\|w\|_{L^2}^2 .
% \end{aligned}
%\end{equation}
\end{proposition}

\begin{proof}
By \eqref{s4:thetabound}, $\theta_t(v)\le\theta_0$. By Lemma~\ref{s4:LEM:recentre}, \eqref{s4:Taylor} and
Lemma~\ref{s4:LEM:gauss}(i),
\begin{equation}\label{s4:Hessratio}
\begin{aligned}
\H V_t(\ph)[h,h]+\|D_th\|_{L^2}^2
&=\frac{\E_\mu\big[\jb{T_{t,v}u,D_t\jb\nb h}^2\,e^{\RT}\IT\big]}{\E_\mu\big[e^{\RT}\IT\big]}\\
&\quad-\bigg(\frac{\E_\mu\big[\jb{T_{t,v}u,D_t\jb\nb h}\,e^{\RT}\IT\big]}{\E_\mu\big[e^{\RT}\IT\big]}\bigg)^2 .
\end{aligned}
\end{equation}
Below, $C$ and the implicit constants in $\les$ depend only on $p$, $K$ and $\theta_0$.

First of all, note that $\g_0 \ll 1$ implies $t \ll 1$.
Since $1-e^{-s}\sim\min(s,1)$ for $s\ge0$, we have
$\sum_{n\in\Z}\frac{1-e^{-t\jb n^2}}{\jb n^2}\sim\sum_{n\in\Z}\min(t,\jb n^{-2})\sim t^{\frac12}$.
So Lemma~\ref{s4:LEM:gauss}(iv) shows that, for every $x\in\T$, $D_t^{-1}\jb\nb T_{t,v}u(x)$ is a centred Gaussian
vector in $\R^2$ with $\E_\mu|D_t^{-1}\jb\nb T_{t,v}u(x)|^2\les t^\frac12$. By \eqref{s4:Taylor} and Minkowski's
inequality, since $t\le1$,
\[
\|\RT\|_{L^4(\mu)}\les t^{\frac34}\int_\T|v|^{p-3}+t^{\frac p4}.
\]
Since $\|v\|_{L^2}^2\le\frac K2$, Chebyshev's inequality gives
\[
\mu\big(\|v+D_t^{-1}\jb\nb T_{t,v}u\|_{L^2}^2>K\big)\les\E_\mu\|D_t^{-1}\jb\nb T_{t,v}u\|_{L^2}^2\les t^\frac12 .
\]

\noi
 By Lemma~\ref{s4:LEM:BD} with $q=4$,
\[
\E_\mu\big[e^{4|\RT|}\IT\big]\les1 .
\]
We write
\begin{align*}
&\mu\big(\|v+D_t^{-1}\jb\nb T_{t,v}u\|_{L^2}^2\le K\big) \\
&= \E_\mu\big[e^{\frac45\RT} e^{-\frac45\RT}\IT\big].
\end{align*}
Therefore, by H\"older we have that
\begin{align*}
\E_\mu\big[e^{\RT}\IT\big]& \gtrsim \mu\big(\|v+D_t^{-1}\jb\nb T_{t,v}u\|_{L^2}^2\le K\big)^\frac54\\
&\gtrsim 1 - \frac 54 \mu\big(\|v+D_t^{-1}\jb\nb T_{t,v}u\|_{L^2}^2>K\big) \gtrsim 1
\end{align*}
for $t \ll 1$.

Now let
\[
\eps:=\bigg\|\frac{e^{\RT}\IT}{\E_\mu\big[e^{\RT}\IT\big]}-1\bigg\|_{L^2(\mu)} .
\]
Since the denominator is bounded below,  the elementary inequality $|e^x-1|\le|x|(e^x+1)$ and H\"older's inequality give
\[
\begin{aligned}
\eps&\les\big\|e^{\RT}\IT-1\big\|_{L^2(\mu)}\\
&\le\|\RT\|_{L^4(\mu)}\Big(\E_\mu\big[e^{4\RT}\IT\big]^{\frac14}+1\Big)\\
&\quad+\mu\big(\|v+D_t^{-1}\jb\nb T_{t,v}u\|_{L^2}^2>K\big)^{\frac12}\les\gamma .
\end{aligned}
\]
So $\eps\les\gamma$, in particular, $\eps\ll 1$.

Write  $Y = \jb{T_{t,v}u,D_t\jb\nb h},$ and note that this is a centred Gaussian under $\mu$. We have that, for every density $\rho$ with $\E[\rho]=1,$
\begin{align*}
    &\big|\E[Y^2 \rho] - \E[Y\rho]^2 - \E[Y^2]\big| \\
    &= \big|\E[Y^2 (\rho-1)] - \E[Y(\rho-1)]^2\big| \\
    &\le \E[Y^4]^\frac12 \|\rho-1\|_{L^2} + \E[Y^2] \|\rho-1\|_{L^2}^2 \\
    &\les \E[Y^2]\|\rho-1\|_{L^2}(1 + \|\rho-1\|_{L^2}),
\end{align*}
since $\E[Y^4]=3\E[Y^2]^2$.
Applying this to
$$ \rho = \frac{e^{\RT}\IT}{\E_\mu\big[e^{\RT}\IT\big]},$$
and noting that the RHS of \eqref{s4:Hessratio} is exactly $\E[Y^2 \rho] - \E[Y\rho]^2$,
this gives
\[
\Big|\H V_t(\ph)[h,h]+\|D_th\|_{L^2}^2-\E_\mu\big[\jb{T_{t,v}u,D_t\jb\nb h}^2\big]\Big|
\les\eps\,\E_\mu\big[\jb{T_{t,v}u,D_t\jb\nb h}^2\big].
\]
By Lemma~\ref{s4:LEM:gauss}(v), $\E_\mu\big[\jb{T_{t,v}u,D_t\jb\nb h}^2\big]\le(1-\theta_0)^{-1}\|D_th\|_{L^2}^2$,
 and \eqref{s4:crit} follows, since $\eps\les\gamma$.
%Finally, \eqref{s4:critw} follows from \eqref{s4:crit} with
% $h=e^{-\frac t2\jb\nb^2}w$, since the symbol $\jb n^2(e^{t\jb n^2}-1)^{-1}$ of $D_t^2e^{-t\jb\nb^2}$ is at most $t^{-1}$,
% so that $\|D_te^{-\frac t2\jb\nb^2}w\|_{L^2}^2\le t^{-1}\|w\|_{L^2}^2$.
\end{proof}

\end{document}
